% ==========================================================================
% Preambulo - estilo clasico institucional UAX
% ==========================================================================
\usepackage[utf8]{inputenc}
\usepackage[T1]{fontenc}
\usepackage[spanish,es-tabla]{babel}
\usepackage{lmodern}

% Nota: LaTeX puede emitir el warning 'Font shape T1/lmr/bx/sc undefined'
% al procesar el TOC. La sustitucion automatica a bold-normal es correcta.

% -------- Geometria --------
\usepackage[a4paper, top=3cm, bottom=3cm, left=3cm, right=2.5cm,
            headheight=14pt, headsep=0.5cm, footskip=1.2cm]{geometry}
\usepackage{setspace}
% Interlineado 1,15 seg\'un normativa UAX
\setstretch{1.15}

% Permitir mas margen de estiramiento en lineas dificiles.
\emergencystretch=3em
\tolerance=1000
\hbadness=10000

% -------- Cabecera y pie clasicos --------
\usepackage{fancyhdr}
\pagestyle{fancy}
\fancyhf{}
\renewcommand{\chaptermark}[1]{\markboth{\thechapter.\ #1}{}}
\renewcommand{\sectionmark}[1]{}
\fancyhead[LE,RO]{\small\thepage}
\fancyhead[LO,RE]{\small\itshape\nouppercase{\leftmark}}
\renewcommand{\headrulewidth}{0.4pt}
\renewcommand{\footrulewidth}{0pt}

\fancypagestyle{plain}{%
  \fancyhf{}
  \fancyfoot[C]{\small\thepage}
  \renewcommand{\headrulewidth}{0pt}
}
\fancypagestyle{frontmatter}{%
  \fancyhf{}
  \fancyfoot[C]{\small\thepage}
  \renewcommand{\headrulewidth}{0pt}
}

% -------- Figuras, tablas, codigo --------
\usepackage{graphicx}
\usepackage{float}
\usepackage{booktabs}
\usepackage{longtable}
\usepackage{multirow}
\usepackage{array}
\usepackage{caption}
\usepackage{subcaption}
\usepackage{xcolor}
\usepackage{listings}

\captionsetup{font=small, labelfont=bf, labelsep=period,
              justification=raggedright, singlelinecheck=false}

\definecolor{codegray}{gray}{0.96}
\definecolor{codestring}{RGB}{163,21,21}
\definecolor{codecomment}{RGB}{92,122,72}
\definecolor{codekw}{RGB}{45,110,175}
\lstset{
  basicstyle=\ttfamily\footnotesize,
  backgroundcolor=\color{codegray},
  frame=single,
  rulecolor=\color{lightgray},
  breaklines=true,
  numbers=left,
  numberstyle=\tiny\color{gray},
  keywordstyle=\color{codekw}\bfseries,
  stringstyle=\color{codestring},
  commentstyle=\color{codecomment}\itshape,
  showstringspaces=false,
  captionpos=b,
  extendedchars=true,
  literate={á}{{\'a}}1 {é}{{\'e}}1 {í}{{\'i}}1 {ó}{{\'o}}1 {ú}{{\'u}}1
           {Á}{{\'A}}1 {É}{{\'E}}1 {Í}{{\'I}}1 {Ó}{{\'O}}1 {Ú}{{\'U}}1
           {ñ}{{\~n}}1 {Ñ}{{\~N}}1 {ü}{{\"u}}1 {Ü}{{\"U}}1
}
\lstdefinestyle{bash}{language=bash}
\lstdefinestyle{python}{language=Python}

% -------- Listas con opciones --------
\usepackage{enumitem}

% -------- Matematicas --------
\usepackage{amsmath,amssymb,mathtools}

% -------- Referencias cruzadas y bibliografia --------
% Azul institucional UAX para todos los enlaces (interno y externo)
\definecolor{uaxblue}{HTML}{1866B6}

\usepackage{hyperref}
\hypersetup{
  colorlinks=true,
  linkcolor=uaxblue,
  citecolor=uaxblue,
  urlcolor=uaxblue,
  filecolor=uaxblue,
  linktoc=all,
  pdftitle={Framework transcriptomico para biomarcadores en cancer de pulmon},
  pdfauthor={Daniel Tapia Diez},
  pdfsubject={TFM Bioinformatica, UAX},
  pdfkeywords={transcriptomica, cancer de pulmon, machine learning, biomarcadores, LODO, LLM, NCBI GEO}
}
\usepackage[numbers,sort&compress]{natbib}
\bibliographystyle{unsrtnat}

% Espacio en tablas
\renewcommand{\arraystretch}{1.15}

% -------- Titulos clasicos: "Capitulo N" arriba, titulo Huge debajo --------
\usepackage{titlesec}

\titleformat{\chapter}[display]
  {\normalfont\huge\bfseries}
  {\chaptertitlename\ \thechapter}
  {20pt}
  {\Huge}
\titlespacing*{\chapter}{0pt}{-20pt}{20pt}

\titleformat{\section}
  {\normalfont\Large\bfseries}
  {\thesection.}{0.6em}{}
\titlespacing*{\section}{0pt}{1.2em}{0.6em}

\titleformat{\subsection}
  {\normalfont\large\bfseries}
  {\thesubsection.}{0.5em}{}
\titlespacing*{\subsection}{0pt}{1em}{0.4em}

\titleformat{\subsubsection}
  {\normalfont\normalsize\bfseries}
  {\thesubsubsection.}{0.5em}{}

% -------- Comandos utiles (SIN color verde) --------
\newcommand{\framework}{\texttt{data.lung}}
\newcommand{\gene}[1]{\textit{#1}}
\newcommand{\cohorte}[1]{\texttt{#1}}

% -------- Anexos --------
\usepackage[toc,page]{appendix}
